\begin{textsample}{POS Dim 1 – se – Score 29.00 – se/se\_ef\_2.txt}  \label{ex:f1_pos_120}
ARTE A Educação \textbf{atual} vem trabalhando num percurso onde o grande desafio é criar \textbf{uma} proposta pedagógica capaz de acompanhar as necessidades do ser humano \textbf{moderno}, \textbf{que} faz parte de \textbf{uma} sociedade \textbf{globalizada} \textbf{que} \textbf{exige} \textbf{uma} educação mais abrangente, contextualizada e integral.
No ensino da Arte, o caminho não é diferente : o sentir e o pensar traduzem, nitidamente, o \textbf{que} a mesma propõe : trabalhar na criança e no adolescente a sensibilidade \textbf{inerente} ao ser humano.
Na visão de Freitas ( 2011, s/p ), a escola \textbf{atual} deve fazer com \textbf{que} seu \textbf{alunado} pense, reflita, analise, sintetize, critique, crie, classifique, tire conclusões, estabeleça relações, \textbf{argumente}, avalie, justifique etc.
Para \textbf{que} todo este processo aconteça, é de fundamental importância \textbf{que} os professores trabalhem com metodologias participativas, desafiadoras, problematizando os conteúdos e levando o aluno a pensar, a formular hipóteses, a descobrir, a falar, a questionar, a colocar suas opiniões, suas divergências e dúvidas, a trocar informações com o grupo de colegas, defendendo e \textbf{argumentando} seu ponto de vista.
Conforme Libâneo, Oliveira e Toschi ( 2009, p. 994 ), “ A escola é \textbf{uma} organização em \textbf{que} tanto seus objetivos e resultados quanto seus processos e meios são relacionados com a formação humana, ganhando relevância, portanto, o fortalecimento das relações sociais, culturais e afetivas \textbf{que} nela têm lugar. ” Especificamente no campo das Artes, \textbf{hoje} a BNCC mais \textbf{uma} vez contempla esta \textbf{disciplina}, entendendo a real importância \textbf{que} a mesma tem dentro do contexto histórico, cultural e social da \textbf{humanidade}.
Afinal, foi através da Arte \textbf{que} perpetuamos e pudemos conhecer toda a trajetória humana ao longo da história.
A Arte interage com todas as demais \textbf{disciplinas} e oportuniza ao \textbf{alunado} a ter \textbf{uma} visão mais ampla sobre sua história e sua própria vida, pois a capacidade de criar é intrínseca ao \textbf{homem}.
David Ausubel ( 2001, p. 76 ) traz \textbf{uma} \textbf{teoria} em relação à aprendizagem e esta vem ratificar o \textbf{que} queremos a partir da reelaboração de nosso currículo de Artes.
Para ele, “ [... ] a aprendizagem será muito mais significativa na medida em \textbf{que} novo material for incorporado às estruturas de conhecimento de \textbf{um} aluno e este adquira significado para o mesmo a partir da relação com o seu conhecimento prévio. ” Desta forma, é de \textbf{suma} importância trazer a realidade do aluno para dentro da sala de aula.
Não obstante, não sendo feito assim, a aprendizagem tornar -se-á mecânica ou repetitiva, \textbf{visto} \textbf{que} vai se produzir menos essa incorporação e atribuição do significado ; assim, o \textbf{que} se apresenta como novo para o aluno, sem sua identificação pessoal, apenas vai ser armazenado isoladamente ou por meio de associações arbitrárias na estrutura cognitiva.
Infelizmente, desde muito, a escola tradicional vem enfatizando as habilidades manuais, os “ dons ” em detrimento de \textbf{um} saber mais amplo de Arte, \textbf{que} valorizasse a criação pessoal.
Na prática, o ensino de Arte estava mais direcionado ao domínio técnico formador da qualificação profissional para o mercado de trabalho, especialmente para as classes menos favorecidas.
A classe mais abastada poderia desenvolver o gosto apreciativo pela arte.
Percebe -se, nitidamente, o poder segregador \textbf{que} a \textbf{disciplina} Arte carregava consigo.
De acordo com Marques ( 2001, p. 32 ), “ somente no final da década de 1990, entidades, associações e órgãos governamentais \textbf{preocuparam} -se em incluir outras linguagens artísticas nas discussões acerca do ensino de Arte ”.
A BNCC vem trazer para a Educação Básica, incluindo a \textbf{disciplina} Arte, \textbf{uma} visão centrada na formação integral de nosso \textbf{alunado}, desmistificando a função segregadora assumida por muito tempo na história da Educação Brasileira.
O presente instrumento, \textbf{que} a partir desta reformulação, passa a se denominar Currículo Básico Comum, é fruto das pesquisas \textbf{que} temos realizado através da BNCC e das capacitações das quais participamos, \textbf{que} nos permitiram o contato com diversos professores do Brasil e do próprio Estado, experimentando várias discussões acerca do aluno \textbf{que} se quer formar a partir de agora.
Com este trabalho, vamos oferecer o primeiro currículo, de fato, do estado de Sergipe, \textbf{visto} \textbf{que} antes, apenas tínhamos \textbf{um} Referencial Curricular \textbf{que}, como o nome já ratifica, era apenas \textbf{um} referencial para os professores e para a escola na construção dos PPPs, dos planejamentos anuais e planos de aula consecutivamente.
Neste momento, nosso objetivo é trazer para a \textbf{disciplina} Arte \textbf{uma} proposta mais dinâmica de aprendizagem, com a preocupação de se fazer \textbf{um} currículo mais acessível a todos.
Intencionalmente, trabalhamos \textbf{um} currículo genuinamente sergipano, \textbf{que} trabalhe as especificidades de nosso Estado, de nossos municípios e, através dele, nosso \textbf{alunado} sinta -se verdadeiramente sergipano, conhecendo mais sua cultura, seu folclore, suas tradições e perceba -se \textbf{um} cidadão pró-ativo no lugar onde \textbf{vive}.
Além do mais, nossa proposta foi também fazer com \textbf{que} as aulas de Arte possam oportunizar a nosso \textbf{educando} trazer para a sala de aula seu cotidiano, dividindo suas experiências, suas vivências, seus gostos e podendo compartilhar com os demais, construindo, assim, \textbf{uma} arte mais humana e dinâmica, em constante evolução.
Destacamos \textbf{que}, por diversas vezes, sugerimos o trabalho interdisciplinar.
Acreditamos \textbf{que} o trabalho conjunto seja \textbf{uma} metodologia significativa para potencializar o processo de ensino e aprendizagem.
Muitos de nossos conteúdos e habilidades guardam \textbf{interfaces} com os demais componentes curriculares e a construção do trabalho conjunto deve ser \textbf{uma} preocupação permanente de todo o corpo docente da escola.
Na reflexão sobre o ensino de Arte, em qualquer etapa da escolarização, é necessário, como ponto de partida, olharmos de perto o aluno do ano escolar em questão.
Quais são seus interesses, o \textbf{que} já sabe acerca dos fenômenos relacionados aos conteúdos \textbf{que} serão estudados, \textbf{que} tipo de dificuldades apresenta nessa etapa de sua formação, quais são suas expectativas nesse ano escolar.
E, a partir daí, construir com ele os saberes novos, possibilitar o desenvolvimento das habilidades básicas, necessárias ao seu processo de aprendizagem.
O professor poderá trabalhar com vários livros, enciclopédias, com textos e atividades diversas, como observação da Arte construída no entorno, usando a criatividade de acordo com o assunto proposto.
Ele deverá discutir e propor atividades de Arte \textbf{que} poderão ser desenvolvidas em sala de aula, para \textbf{que} os alunos sejam capazes de ler e compreender os gêneros textuais específicos da \textbf{disciplina}, familiarizando -se com a linguagem artística pertinente a cada \textbf{um} dos eixos, isto é, das Artes Visuais, da Dança, da Música, do Teatro e das Artes Integradas, estabelecendo relação entre o \textbf{que} se conhece e o \textbf{que} se lê e produz.
Esses eixos estão representados no organizador curricular de Artes como Unidades Temáticas, \textbf{que} reúnem objetos de conhecimento, especificação dos objetivos de conhecimento e habilidades.
Finalmente, para o desenvolvimento das habilidades, ao longo dos anos de escolaridade, o professor deve mobilizar conhecimentos prévios, contextualizando, \textbf{despertando} a atenção e o apreço do aluno para a temática.
Faz -se necessário, em outro momento, aprofundar essa habilidade, num trabalho \textbf{sistematizado}, relacionando essas aprendizagens ao contexto e a outros temas próximos.
Finalmente, consolidar aquela aprendizagem, também com atividades \textbf{sistematizadas}, para torná -la \textbf{um} saber significativo para o aluno, com o qual ele possa contar para desenvolver outras habilidades, ao longo de seu processo educacional.
Em articulação com as competências gerais da BNCC e as competências específicas da área de Linguagens, o componente curricular de Arte deve garantir aos alunos o desenvolvimento das seguintes competências específicas : 1.
Explorar, conhecer, fruir e analisar criticamente práticas e produções artísticas e culturais do seu entorno social, dos povos indígenas, das comunidades tradicionais brasileiras e de diversas sociedades, em distintos tempos e espaços, para reconhecer a arte como \textbf{um} fenômeno cultural, histórico, social e sensível a diferentes contextos e dialogar com as diversidades. 2.
Compreender as relações entre as linguagens da Arte e suas práticas integradas, inclusive aquelas possibilitadas pelo uso das novas tecnologias de informação e comunicação, pelo cinema e pelo audiovisual, nas condições particulares de produção, na prática de cada linguagem e nas suas articulações. 3.
Pesquisar e conhecer distintas matrizes estéticas e culturais – especialmente aquelas manifestas na arte e nas culturas \textbf{que} constituem a identidade brasileira –, sua tradição e manifestações \textbf{contemporâneas}, reelaborando -as nas criações em Arte. 4.
Experienciar a \textbf{ludicidade}, a \textbf{percepção}, a expressividade e a imaginação, ressignificando espaços da escola e de fora dela no âmbito da Arte. 5.
Mobilizar recursos tecnológicos como formas de registro, pesquisa e criação artística. 6.
Estabelecer relações entre arte, mídia, mercado e consumo, compreendendo, de forma crítica e problematizadora, modos de produção e de circulação da arte na sociedade. 7.
Problematizar questões políticas, sociais, econômicas, científicas, tecnológicas e culturais, por meio de exercícios, produções, intervenções e apresentações artísticas. 8.
Desenvolver a autonomia, a crítica, a autoria e o trabalho coletivo e colaborativo nas artes. 9.
Analisar e valorizar o patrimônio artístico nacional e internacional, material e imaterial, com suas histórias e diferentes visões de mundo.
Os conteúdos foram pensados e estruturados visando à construção de conhecimentos \textbf{que} devem fazer parte da vida de todo ser humano.
Possuem unidade \textbf{conceitual}, \textbf{que} não é seriada e \textbf{que} permite ao professor iniciar o entendimento da Arte a partir de qualquer \textbf{um} dos \textbf{tópicos}.
Possibilita, ainda, a expansão do conhecimento pela criação de redes de informação.

% matched lemmas: alunado, argumentar, atual, conceitual, contemporâneo, despertar, disciplina, educar, exigir, globalizado, hoje, homem, humanidade, inerente, interface, ludicidade, moderno, percepção, preocupar, que, sistematizar, sumo, teoria, tópico, um, ver, viver
\end{textsample}
